% !TEX TS-program = pdflatex
% !TEX encoding = UTF-8 Unicode

% This file is a template using the "beamer" package to create slides for a talk or presentation
% - Giving a talk on some subject.
% - The talk is between 15min and 45min long.
% - Style is ornate.

% MODIFIED by Jonathan Kew, 2008-07-06
% The header comments and encoding in this file were modified for inclusion with TeXworks.
% The content is otherwise unchanged from the original distributed with the beamer package.

\documentclass{beamer}


% Copyright 2004 by Till Tantau <tantau@users.sourceforge.net>.
%
% In principle, this file can be redistributed and/or modified under
% the terms of the GNU General Public License, version 2.
%
% However, this file is supposed to be a template to be modified
% for your own needs. For this reason, if you use this file as a
% template and not specifically distribute it as part of a another
% package/program, I grant the extra permission to freely copy and
% modify this file as you see fit and even to delete this copyright
% notice. 


\mode<presentation>
{
  \usetheme{Warsaw}
  % or ...

  \setbeamercovered{transparent}
  % or whatever (possibly just delete it)
}


\usepackage[english]{babel}
\usepackage[utf8]{inputenc}
\usepackage[]{hyperref}

\usepackage{times}
\usepackage[T1]{fontenc}
% Or whatever. Note that the encoding and the font should match. If T1
% does not look nice, try deleting the line with the fontenc.

\begin{document}

\title[Short Paper Title] % (optional, use only with long paper titles)
{Presentation Title}

\subtitle
{Presentation Subtitle} % (optional)

\author[Author, Another] % (optional, use only with lots of authors)
{F.~Author\inst{1} \and S.~Another\inst{2}}
% - Use the \inst{?} command only if the authors have different
%   affiliation.

\institute[Universities of Somewhere and Elsewhere] % (optional, but mostly needed)
{
  \inst{1}%
  Department of Computer Science\\
  University of Somewhere
  \and
  \inst{2}%
  Department of Theoretical Philosophy\\
  University of Elsewhere}
% - Use the \inst command only if there are several affiliations.
% - Keep it simple, no one is interested in your street address.

\date[Short Occasion] % (optional)
{Date / Occasion}

\subject{Talks}
%
\begin{frame}{Launch System Hierarchy}
    \centering
    \footnotesize
    
    % Set up a customized table
    % - First column (r): Right-aligns the numbering to line up the rightmost decimals
    % - @{\hspace{0.3cm}}: Forces a clean, consistent structural indent gap
    % - Second column (l): Left-aligns your system text entries
    
    \begin{tabular}{ l @{\hspace{0.3cm}} l }
        \textbf{WBS} & \textbf{Product Element} \\
         1.0   & Strategic Missile Weapon System \\
        1.2.1 & Launch Facility Equipment \\
        1.2.2 & Launch Support Equipment \\
        1.2.3 & Launch Control Interfaces \\
        1.2.4 & Environmental Control Equipment \\
    \end{tabular}

\end{frame}
%
\begin{frame}{Launch System Hierarchy}
    \centering
    \footnotesize
    
    % Set up a customized table
    % - First column (r): Right-aligns the numbering to line up the rightmost decimals
    % - @{\hspace{0.3cm}}: Forces a clean, consistent structural indent gap
    % - Second column (l): Left-aligns your system text entries
    
    \begin{tabular}{ l @{\hspace{0.3cm}} l }
        \textbf{WBS} & \textbf{Product Element} \\
         1.2   & Launch System \\
        1.2.1 & Launch Facility Equipment \\
        1.2.2 & Launch Support Equipment \\
        1.2.3 & Launch Control Interfaces \\
        1.2.4 & Environmental Control Equipment \\
    \end{tabular}
\end{frame}

\section*{Summary}

\begin{frame}{Summary}

  % Keep the summary *very short*.
  \begin{itemize}
  \item
    The \alert{first main message} of your talk in one or two lines.
  \item
    The \alert{second main message} of your talk in one or two lines.
  \item
    Perhaps a \alert{third message}, but not more than that.
  \end{itemize}
  
  % The following outlook is optional.
  \vskip0pt plus.5fill
  \begin{itemize}
  \item
    Outlook
    \begin{itemize}
    \item
      Something you haven't solved.
    \item
      Something else you haven't solved.
    \end{itemize}
  \end{itemize}
\end{frame}


\end{document}


